\documentclass[12pt, a4paper]{article}

\usepackage[utf8]{inputenc}
\usepackage[T1]{fontenc}
\usepackage[english]{babel}
\usepackage{geometry}
\usepackage{amsmath}
\usepackage{amssymb}
\usepackage{booktabs}
\usepackage{hyperref}
\usepackage{xcolor}
\usepackage{listings}
\usepackage{graphicx}
\usepackage{float}
\usepackage{fancyhdr}
\usepackage{parskip}

\geometry{top=2.5cm, bottom=2.5cm, left=3cm, right=2.5cm, headheight=15pt}

\hypersetup{
    colorlinks=true,
    linkcolor=blue!70!black,
    urlcolor=blue!70!black,
    citecolor=blue!70!black
}

\definecolor{codegray}{rgb}{0.95,0.95,0.95}
\definecolor{codecomment}{rgb}{0.4,0.4,0.4}
\definecolor{codestring}{rgb}{0.13,0.54,0.13}
\definecolor{codekeyword}{rgb}{0.0,0.0,0.8}

\lstset{
    backgroundcolor=\color{codegray},
    commentstyle=\color{codecomment}\itshape,
    keywordstyle=\color{codekeyword}\bfseries,
    stringstyle=\color{codestring},
    basicstyle=\ttfamily\small,
    breakatwhitespace=false,
    breaklines=true,
    keepspaces=true,
    numbers=left,
    numbersep=8pt,
    numberstyle=\tiny\color{codecomment},
    showspaces=false,
    showstringspaces=false,
    showtabs=false,
    tabsize=4,
    frame=single,
    framerule=0.4pt,
    rulecolor=\color{gray!40}
}

\pagestyle{fancy}
\fancyhf{}
\fancyhead[L]{\small Home Credit Default Risk}
\fancyhead[R]{\small Valvitor Santos}
\fancyfoot[C]{\thepage}
\renewcommand{\headrulewidth}{0.4pt}

\title{
    \vspace{-1cm}
    {\Large \textbf{Home Credit Default Risk}} \\[0.5em]
    {\large Technical Report --- End-to-End Machine Learning Pipeline\\
    for Loan Default Prediction}
}
\author{Valvitor Santos}
\date{\today}

% ──────────────────────────────────────────────────────────────
\begin{document}

\maketitle
\thispagestyle{fancy}

\begin{center}
\begin{minipage}{0.85\textwidth}
\small\textbf{Abstract.}
This report describes the development of a complete machine learning pipeline for the
\textit{Home Credit Default Risk} Kaggle competition. The problem consists of predicting
the probability of loan default for clients who may have limited or no credit history.
The solution employs feature engineering across seven relational tables, a LightGBM model
with Optuna hyperparameter optimization, and 5-fold stratified cross-validation.
Experiments are tracked with MLflow. The evaluation metric is \textbf{AUC-ROC}.
\end{minipage}
\end{center}

\bigskip
\tableofcontents
\newpage

% ══════════════════════════════════════════════════════════════
\section{Introduction and Problem Context}

\subsection{Motivation}

Financial institutions face the challenge of assessing the repayment capacity of clients
who have little or no formal credit history. The absence of historical data in conventional
credit bureaus leads to the financial exclusion of populations that are, in fact, capable
of honoring their commitments. Home Credit, an organization specialized in responsible
lending, seeks to use alternative data sources to make this assessment fairer and more
accurate.

\subsection{Problem Formulation}

The objective is to build a \textbf{binary classification} model that, given an
applicant's profile, estimates the probability of default:

\begin{equation}
    \hat{p}_i = P(\text{TARGET}_i = 1 \mid \mathbf{x}_i)
\end{equation}

where $\text{TARGET}_i = 1$ indicates that client $i$ did not repay the loan, and
$\mathbf{x}_i$ is the client's feature vector.

The official evaluation metric is the \textbf{AUC-ROC} (Area Under the ROC Curve),
which measures the model's ability to separate defaulters from non-defaulters regardless
of the classification threshold:

\begin{equation}
    \text{AUC} = \int_0^1 \text{TPR}\!\left(\text{FPR}^{-1}(t)\right) dt
\end{equation}

\subsection{Main Challenges}

\begin{itemize}
    \item \textbf{Class imbalance:} approximately 8\% of contracts result in default,
          creating strong bias toward the majority class.
    \item \textbf{Data distributed across multiple tables:} relevant information is
          fragmented across 7 relational tables that must be aggregated to the applicant
          level.
    \item \textbf{High proportion of missing values:} some features have up to 60\%
          null values, requiring careful treatment.
    \item \textbf{Critical feature engineering:} model performance depends heavily on
          the quality of derived features, especially from auxiliary tables.
\end{itemize}

% ══════════════════════════════════════════════════════════════
\section{Dataset}

\subsection{Relational Structure}

The data is provided in 7 CSV files with the following hierarchy:

\begin{table}[H]
\centering
\begin{tabular}{llrl}
\toprule
\textbf{Table} & \textbf{Key} & \textbf{Rows} & \textbf{Description} \\
\midrule
\texttt{application\_train/test} & \texttt{SK\_ID\_CURR}   & 307,511 / 48,744 & Current loan application \\
\texttt{bureau}                  & \texttt{SK\_ID\_CURR}   & 1,716,428        & Credit at other banks \\
\texttt{bureau\_balance}         & \texttt{SK\_ID\_BUREAU} & 27,299,925       & Monthly bureau history \\
\texttt{previous\_application}   & \texttt{SK\_ID\_CURR}   & 1,670,214        & Prior HC applications \\
\texttt{POS\_CASH\_balance}      & \texttt{SK\_ID\_CURR}   & 10,001,358       & POS and cash loan balances \\
\texttt{installments\_payments}  & \texttt{SK\_ID\_CURR}   & 13,605,401       & Repayment history \\
\texttt{credit\_card\_balance}   & \texttt{SK\_ID\_CURR}   & 3,840,312        & Credit card balances \\
\bottomrule
\end{tabular}
\caption{Dataset structure}
\end{table}

The main table is the model anchor; all others are aggregated to the \texttt{SK\_ID\_CURR}
level and joined via \texttt{left join}.

\subsection{Target Variable}

\begin{itemize}
    \item $\texttt{TARGET} = 0$: client repaid normally (91.93\%)
    \item $\texttt{TARGET} = 1$: client defaulted (8.07\%)
\end{itemize}

% ══════════════════════════════════════════════════════════════
\section{Exploratory Data Analysis}

The EDA was conducted in notebook \texttt{01\_EDA.ipynb} with the goal of extracting
actionable insights --- not merely descriptive statistics, but pre-processing and feature
engineering decisions directly informed by the data.

\subsection{Anomaly in \texttt{DAYS\_EMPLOYED}}

The variable \texttt{DAYS\_EMPLOYED} represents how many days a client has been employed
(negative value). However, 18,148 clients show the value \textbf{365,243} --- clearly a
placeholder for ``not employed'', equivalent to $\approx$ 1,000 years of employment.

\textbf{Decision:} replace 365,243 with \texttt{NaN} and create a binary feature
\texttt{FLAG\_EMPLOYED\_ANOMALY} to preserve the information that data was missing.

\begin{lstlisting}[language=Python, caption={Handling the DAYS\_EMPLOYED anomaly}]
df['FLAG_EMPLOYED_ANOMALY'] = (df['DAYS_EMPLOYED'] == 365243).astype(np.int8)
df['DAYS_EMPLOYED'] = df['DAYS_EMPLOYED'].replace(365243, np.nan)
\end{lstlisting}

\subsection{EXT\_SOURCE Features}

The variables \texttt{EXT\_SOURCE\_1}, \texttt{EXT\_SOURCE\_2}, and
\texttt{EXT\_SOURCE\_3} represent external credit scores from third-party bureaus. EDA
revealed these are the \textbf{most correlated features with default} in the entire
dataset, making them priority candidates for non-linear combinations.

\texttt{EXT\_SOURCE\_1} has 56\% missing values, requiring an additional flag:
\texttt{FLAG\_EXT\_SOURCE\_1\_MISSING}.

\subsection{Real Estate Block}

Approximately 20 variables related to the client's property (\texttt{APARTMENTS\_AVG},
\texttt{BASEMENTAREA\_AVG}, etc.) have more than 60\% null values and low correlation
with the target.

\textbf{Decision:} drop all columns with more than 60\% missing values, except
\texttt{EXT\_SOURCE}.

\subsection{Class Imbalance}

With 8.07\% default rate, unadjusted models tend to classify almost all clients as
non-defaulters. \textbf{Decision:} use LightGBM's \texttt{scale\_pos\_weight} parameter:

\begin{equation}
    \texttt{scale\_pos\_weight} = \frac{N_{\text{negatives}}}{N_{\text{positives}}}
    = \frac{283{,}301}{24{,}210} \approx 11.70
\end{equation}

% ══════════════════════════════════════════════════════════════
\section{Feature Engineering}

\subsection{Overview}

Feature engineering was executed in notebook \texttt{02\_feature\_engineering.ipynb},
transforming 7 raw tables into a single processed dataset with over 1,000 features.
The process follows four steps:

\begin{enumerate}
    \item Manual features on the main application table
    \item Bureau and bureau\_balance aggregation
    \item Auxiliary table aggregation (previous, POS, installments, credit card)
    \item Final merge, One-Hot Encoding and alignment
\end{enumerate}

\subsection{Main Table Features}

\subsubsection{Financial Ratios}

Ratios capturing the relationship between credit, income, and installments:

\begin{table}[H]
\centering
\begin{tabular}{ll}
\toprule
\textbf{Feature} & \textbf{Formula} \\
\midrule
\texttt{CREDIT\_INCOME\_RATIO}  & $\texttt{AMT\_CREDIT} / (\texttt{AMT\_INCOME\_TOTAL} + 1)$ \\
\texttt{ANNUITY\_INCOME\_RATIO} & $\texttt{AMT\_ANNUITY} / (\texttt{AMT\_INCOME\_TOTAL} + 1)$ \\
\texttt{ANNUITY\_CREDIT\_RATIO} & $\texttt{AMT\_ANNUITY} / (\texttt{AMT\_CREDIT} + 1)$ \\
\texttt{INCOME\_PER\_PERSON}    & $\texttt{AMT\_INCOME\_TOTAL} / (\texttt{CNT\_FAM\_MEMBERS} + 1)$ \\
\texttt{GOODS\_CREDIT\_RATIO}   & $\texttt{AMT\_GOODS\_PRICE} / (\texttt{AMT\_CREDIT} + 1)$ \\
\bottomrule
\end{tabular}
\caption{Financial features created from the main table}
\end{table}

Adding 1 to the denominator avoids division by zero without distorting values for
large denominators.

\subsubsection{Time-Based Features}

\begin{itemize}
    \item \texttt{DAYS\_BIRTH\_YEARS}: client age in years (more interpretable than
          negative days)
    \item \texttt{DAYS\_EMPLOYED\_YEARS}: employment duration in years
    \item \texttt{EMPLOYED\_TO\_BIRTH\_RATIO}: fraction of life spent working
    \item \texttt{EMPLOYED\_TO\_AGE\_PCT}: employment duration / age ratio
\end{itemize}

\subsubsection{EXT\_SOURCE Combinations}

Given the high predictive power of external scores, multiple combinations were created:

\begin{align}
    \texttt{EXT\_SOURCE\_MEAN} &= \tfrac{1}{3}(E_1 + E_2 + E_3) \\
    \texttt{EXT\_SOURCE\_STD}  &= \text{std}(E_1, E_2, E_3) \\
    \texttt{EXT\_SOURCE\_PROD} &= E_1 \cdot E_2 \cdot E_3 \\
    \texttt{EXT12} &= E_1 \cdot E_2, \quad
    \texttt{EXT13} = E_1 \cdot E_3, \quad
    \texttt{EXT23} = E_2 \cdot E_3
\end{align}

The product captures non-linear interactions: a low score from any single source
collapses the product, signaling concentrated risk.

\subsection{Auxiliary Table Aggregation}

For each auxiliary table, the process is:

\begin{enumerate}
    \item Create domain-specific manual features for the table
    \item Apply statistical aggregations for numeric columns:
          \texttt{count}, \texttt{mean}, \texttt{max}, \texttt{min}, \texttt{sum}
    \item Apply One-Hot Encoding + aggregation for categorical columns:
          \texttt{sum}, \texttt{mean}
    \item Left-join to the \texttt{SK\_ID\_CURR} level
\end{enumerate}

This process is implemented in reusable functions \texttt{agg\_numeric()} and
\texttt{agg\_categorical()} in \texttt{src/features.py}.

\subsubsection{Bureau and Bureau Balance}

\texttt{bureau\_balance} provides monthly history for each bureau credit entry.
Status columns were binarized:

\begin{lstlisting}[language=Python, caption={Feature engineering on bureau\_balance}]
# C = closed, X = no loan, 1-5 = months overdue (increasing severity)
bureau_balance['STATUS_GOOD'] = bureau_balance['STATUS'].isin(['C','X']).astype(np.int8)
bureau_balance['STATUS_LATE'] = bureau_balance['STATUS'].isin(['1','2','3','4','5']).astype(np.int8)
\end{lstlisting}

\subsubsection{Installments Payments}

\begin{itemize}
    \item \texttt{PAYMENT\_DIFF}: amount paid $-$ amount due (positive = paid more
          than the minimum)
    \item \texttt{PAYMENT\_RATIO}: ratio of amount paid to amount due
    \item \texttt{DAYS\_PAST\_DUE}: days overdue (clipped to 0 for early payments)
    \item \texttt{PAID\_ON\_TIME}: binary on-time payment flag
\end{itemize}

\subsubsection{Credit Card Balance}

\begin{equation}
    \texttt{UTILIZATION} = \frac{\texttt{AMT\_BALANCE}}{\texttt{AMT\_CREDIT\_LIMIT\_ACTUAL} + 1}
\end{equation}

High credit utilization is a strong indicator of financial stress.

\subsubsection{Previous Applications}

\begin{itemize}
    \item \texttt{APP\_CREDIT\_RATIO}: ratio of amount applied for to amount approved
          --- captures whether clients typically request more than they receive
    \item \texttt{DAYS\_DECISION\_YEARS}: time since last credit decision, in years
\end{itemize}

\subsection{Final Processing}

After merging all tables:

\begin{enumerate}
    \item Drop columns with \textbf{more than 80\% nulls} post-merge (auxiliary tables
          introduce many nulls for clients without history in a given product)
    \item \textbf{One-Hot Encoding} for categorical variables with fewer than 10 unique
          categories
    \item \textbf{Alignment} between train and test: columns missing in test are
          filled with 0
    \item \textbf{Column name sanitization}: special characters replaced by \texttt{\_}
          (required for LightGBM compatibility)
    \item Save to \textbf{Parquet format} for efficient I/O
\end{enumerate}

The final result is a dataset with \textbf{1,042 features} and 307,511 training rows.

% ══════════════════════════════════════════════════════════════
\section{Modeling}

\subsection{Algorithm Selection}

\textbf{LightGBM} (Light Gradient Boosting Machine) was chosen for the following reasons:

\begin{itemize}
    \item \textbf{Efficiency with high dimensionality:} the GOSS and EFB algorithms
          allow fast training with over 1,000 features
    \item \textbf{Native missing value handling:} learns the best split direction for
          missing data without imputation
    \item \textbf{Performance on tabular data:} gradient boosting trees historically
          outperform neural networks on tabular data with domain-engineered features
    \item \textbf{Imbalance support:} \texttt{scale\_pos\_weight} adjusts the minority
          class weight in gradient computation
\end{itemize}

\subsection{Hyperparameter Optimization with Optuna}

Optuna implements the \textbf{TPE} (Tree-structured Parzen Estimator), a Bayesian
optimization method that builds a probabilistic model of the hyperparameter space and
suggests configurations with the highest probability of improvement over the current best.

The tuning process used:
\begin{itemize}
    \item 20 trials, each evaluated with 3-fold CV
    \item Objective metric: mean AUC-ROC across the 3 folds
    \item Early stopping of 30 rounds per fold to speed up the search
\end{itemize}

\begin{table}[H]
\centering
\begin{tabular}{lll}
\toprule
\textbf{Hyperparameter} & \textbf{Search Space} & \textbf{Type} \\
\midrule
\texttt{learning\_rate}      & $[0.02,\; 0.15]$   & log-uniform \\
\texttt{num\_leaves}         & $[20,\; 200]$       & integer \\
\texttt{max\_depth}          & $[4,\; 10]$         & integer \\
\texttt{min\_child\_samples} & $[20,\; 300]$       & integer \\
\texttt{feature\_fraction}   & $[0.4,\; 0.9]$      & uniform \\
\texttt{bagging\_fraction}   & $[0.4,\; 0.9]$      & uniform \\
\texttt{bagging\_freq}       & $[1,\; 7]$          & integer \\
\texttt{reg\_alpha}          & $[10^{-4},\; 5.0]$  & log-uniform \\
\texttt{reg\_lambda}         & $[10^{-4},\; 5.0]$  & log-uniform \\
\bottomrule
\end{tabular}
\caption{Optuna search space}
\end{table}

\subsection{Final Training with 5-Fold CV}

After identifying the best hyperparameters, the final model is trained with
\textbf{StratifiedKFold} of 5 folds, which preserves the default rate proportion in
each fold:

\begin{lstlisting}[language=Python, caption={Stratified cross-validation training}]
skf = StratifiedKFold(n_splits=5, shuffle=True, random_state=42)
oof_preds  = np.zeros(len(X))
test_preds = np.zeros(len(X_test))

for fold, (tr_idx, val_idx) in enumerate(skf.split(X, y)):
    model = lgb.LGBMClassifier(**BEST_PARAMS)
    model.fit(X.iloc[tr_idx], y.iloc[tr_idx],
              eval_set=[(X.iloc[val_idx], y.iloc[val_idx])],
              callbacks=[lgb.early_stopping(100)])
    oof_preds[val_idx]  = model.predict_proba(X.iloc[val_idx])[:, 1]
    test_preds         += model.predict_proba(X_test)[:, 1] / 5
\end{lstlisting}

Test predictions are the \textbf{average} across the 5 models, reducing variance.

\subsection{OOF AUC vs.\ Mean Fold AUC}

The \textbf{OOF AUC} (Out-of-Fold) is computed by concatenating all out-of-fold
predictions and calculating AUC over the entire training set at once:

\begin{equation}
    \text{AUC}_{\text{OOF}} = \text{AUC}\!\left(\mathbf{y}_{\text{train}},\;
    \hat{\mathbf{p}}_{\text{OOF}}\right)
\end{equation}

This differs from the mean of per-fold AUCs:

\begin{equation}
    \overline{\text{AUC}}_{\text{folds}} = \frac{1}{K}
    \sum_{k=1}^{K} \text{AUC}\!\left(\mathbf{y}^{(k)}_{\text{val}},\;
    \hat{\mathbf{p}}^{(k)}_{\text{val}}\right)
\end{equation}

The OOF AUC is more reliable because it evaluates the \textbf{global ranking} of all
predictions together. Individual folds may show higher AUC due to statistically more
separable data --- the OOF metric corrects for this bias.

\subsection{Experiment Tracking with MLflow}

All experiments are logged in MLflow:
\begin{itemize}
    \item Best Optuna trial hyperparameters
    \item Final OOF AUC and per-fold AUC
    \item Number of features used
\end{itemize}

This enables reproducible comparison of future runs with new features or different models.

% ══════════════════════════════════════════════════════════════
\section{Engineering Best Practices}

\subsection{Code Structure}

The separation between notebooks and reusable code (\texttt{src/features.py}) avoids
duplication. The functions \texttt{agg\_numeric()} and \texttt{agg\_categorical()} are
generic --- they accept any DataFrame, primary key, and name prefix, making them
applicable to all auxiliary tables without modification.

\subsection{Memory Management}

With datasets of tens of millions of rows, memory usage is critical:

\begin{itemize}
    \item \textbf{Type downcasting:} integers and floats are converted to the smallest
          type that fits their range (\texttt{int8}, \texttt{int16}, \texttt{float32}),
          typically achieving 40--60\% memory reduction
    \item \textbf{Explicit garbage collection:} after each auxiliary table is processed,
          the object is deleted and \texttt{gc.collect()} is called
    \item \textbf{Parquet:} columnar, compressed intermediate storage instead of CSV
\end{itemize}

\subsection{Data Leakage Prevention}

Train-test column alignment is performed via \texttt{pd.DataFrame.align(join=`left')}
to ensure the test set never contains features derived from target information. When
Target Encoding is applied, K-Fold is used to prevent the target from leaking into
the validation fold.

% ══════════════════════════════════════════════════════════════
\section{Conclusion and Future Work}

\subsection{Conclusion}

This project implements a robust end-to-end pipeline for loan default prediction.
The main technical highlights are:

\begin{enumerate}
    \item Decision-driven EDA, not merely descriptive
    \item Feature engineering grounded in financial domain knowledge
    \item Stratified cross-validation to handle class imbalance
    \item Bayesian hyperparameter optimization with Optuna
    \item Full experiment tracking with MLflow
\end{enumerate}

\subsection{Future Work}

\begin{itemize}
    \item \textbf{Stacking/Blending:} combine LightGBM with XGBoost and CatBoost to
          further reduce variance
    \item \textbf{Target Encoding:} apply to high-cardinality variables such as
          \texttt{ORGANIZATION\_TYPE}
    \item \textbf{SHAP values:} feature selection based on permutation importance to
          remove noisy features
    \item \textbf{Pseudo-labeling:} use high-confidence test predictions to augment
          the training set
\end{itemize}

% ══════════════════════════════════════════════════════════════
\appendix
\section{Dependencies}

\begin{table}[H]
\centering
\begin{tabular}{ll}
\toprule
\textbf{Library} & \textbf{Version} \\
\midrule
pandas       & 2.2.2  \\
numpy        & 1.26.4 \\
scikit-learn & 1.4.2  \\
lightgbm     & 4.x    \\
optuna       & 3.6.1  \\
mlflow       & 2.11.3 \\
pyarrow      & 15.0.2 \\
matplotlib   & 3.8.4  \\
seaborn      & 0.13.2 \\
\bottomrule
\end{tabular}
\caption{Main project dependencies}
\end{table}

\end{document}
